% Public-safe visual system for the community template.
% This file contains presentation only: no thesis-specific structure or content.

% TeX Gyre Schola is a portable Schoolbook-style family bundled with current
% TeX Live and MiKTeX distributions.
\setmainfont{texgyreschola-regular.otf}[
  BoldFont=texgyreschola-bold.otf,
  ItalicFont=texgyreschola-italic.otf,
  BoldItalicFont=texgyreschola-bolditalic.otf
]
\setsansfont{texgyreheros-regular.otf}[
  BoldFont=texgyreheros-bold.otf,
  ItalicFont=texgyreheros-italic.otf,
  BoldItalicFont=texgyreheros-bolditalic.otf
]
\setmonofont{lmmono10-regular.otf}[ItalicFont=lmmono10-italic.otf]

\definecolor{ThesisGrey}{gray}{0.55}
\definecolor{ThesisRule}{gray}{0.18}

% A4 layout inspired by the maintained community edition's established memoir
% presentation. Change these values only after checking every output PDF.
\settrimmedsize{297mm}{210mm}{*}
\settypeblocksize{634pt}{*}{*}
\setulmargins{40mm}{*}{*}
\setlrmarginsandblock{26.5mm}{26.5mm}{*}
\setmarginnotes{17pt}{51pt}{\onelineskip}
\setheadfoot{30pt}{2\onelineskip}
\setheaderspaces{*}{2\onelineskip}{*}
\checkandfixthelayout
\OnehalfSpacing
\raggedbottom
\frenchspacing
\setlength{\emergencystretch}{2em}

% Standard hierarchy: chapters, sections and subsections are numbered. Students
% may add or remove headings; the table of contents is generated automatically.
\setsecnumdepth{subsection}
\maxsecnumdepth{subsubsection}
\maxtocdepth{subsection}

% Distinctive main-matter chapter opener: vertical CHAPTER label, large grey box,
% and a right-aligned title. It does not prescribe chapter names or chapter count.
\makeatletter
\newsavebox{\CommunityChapterBox}
\newcommand{\CommunityChapterMarker}[1][4cm]{%
  \sbox{\CommunityChapterBox}{%
    \resizebox{!}{#1}{\fboxsep=1pt%
      \colorbox{ThesisGrey}{\color{white}\thechapter}}}%
  \rotatebox{90}{%
    \resizebox{\heightof{\usebox{\CommunityChapterBox}}+
      \depthof{\usebox{\CommunityChapterBox}}}{!}{\scshape\so\@chapapp}}\quad
  \raisebox{\depthof{\usebox{\CommunityChapterBox}}}{\usebox{\CommunityChapterBox}}%
}
\newcommand{\CommunityChapterNumber}[1][4cm]{%
  \sbox{\CommunityChapterBox}{\CommunityChapterMarker[#1]}%
  \makebox[0pt][c]{\makebox[1cm][r]{\usebox{\CommunityChapterBox}}}%
}

\makechapterstyle{communitymain}{%
  \renewcommand\chapnamefont{\normalfont\Large\scshape\raggedleft\so}
  \renewcommand\chaptitlefont{\normalfont\Large\bfseries\scshape}
  \renewcommand\chapternamenum{}
  \renewcommand\printchaptername{}
  \renewcommand\printchapternum{\null\hfill\CommunityChapterNumber[2.5cm]\par}
  \renewcommand\afterchapternum{\par\vskip\midchapskip}
  \renewcommand\printchaptertitle[1]{%
    \color{ThesisGrey}\chaptitlefont\raggedleft ##1\par}
}

% Front matter uses the same restrained grey language without a numbered marker.
\makechapterstyle{communityfront}{%
  \renewcommand\chapnamefont{\normalfont\Large\scshape}
  \renewcommand\chaptitlefont{\normalfont\Large\bfseries\scshape}
  \renewcommand\printchaptername{}
  \renewcommand\printchapternum{}
  \renewcommand\afterchapternum{}
  \renewcommand\printchaptertitle[1]{%
    \centering\color{ThesisGrey}\chaptitlefont ##1\par}
  \setlength{\beforechapskip}{42pt}
  \setlength{\afterchapskip}{36pt}
}
\makeatother

% Continuation pages carry a restrained running chapter header and rule. Chapter
% opening and front-matter pages retain only the centred page number.
\makepagestyle{communitymain}
\makeoddfoot{communitymain}{}{\thepage}{}
\makeevenfoot{communitymain}{}{\thepage}{}
\makeheadrule{communitymain}{\textwidth}{\normalrulethickness}
\makeoddhead{communitymain}{}{}{\small\textsc{\leftmark}}
\makeevenhead{communitymain}{}{}{\small\textsc{\leftmark}}
\renewcommand*{\chaptermark}[1]{%
  \markboth{\chaptername\space\thechapter.\space #1}{}}

\makepagestyle{communityplain}
\makeoddfoot{communityplain}{}{\thepage}{}
\makeevenfoot{communityplain}{}{\thepage}{}
\makeoddhead{communityplain}{}{}{}
\makeevenhead{communityplain}{}{}{}
\aliaspagestyle{chapter}{communityplain}

% Declaration-page initial. The surrounding statement remains a placeholder and
% must always be replaced with the student's current official wording.
\newcommand{\DeclarationInitial}[1]{%
  \raisebox{-0.50\height}{%
    \fontsize{58}{58}\selectfont\bfseries\textcolor{ThesisGrey}{#1}}%
  \hspace{-0.18em}%
}
